% EEE3096S Practical 1A report, Group 79
% Fields marked [MEASURE: ...] must be filled in during the lab session.

\documentclass[11pt]{article}

\usepackage[a4paper, margin=2.5cm]{geometry}
\usepackage{amsmath}
\usepackage{graphicx}

\title{\textbf{EEE3096S Practical 1A} \\ STM32 GPIO, Timers, and Interrupts}
\author{Tri Le, LXXTRI004 \\ Gill, SHRGIL004}
\date{\today}

\begin{document}

\maketitle

\section{Hardware Orientation}

\begin{enumerate}
    \item GPIO pins connected to the 8 user LEDs: PB0 to PB7. The board maps a full contiguous byte of Port B to the LED cluster: PB0, PB1 and PB2 drive individual LEDs D1 to D3 and PB3 to PB7 drive the remaining five LEDs, each through a 150 $\Omega$ series resistor.
    \item GPIO pins connected to the 4 push buttons: PA0 (SW0), PA1 (SW1), PA2 (SW2) and PA3 (SW3). The switches pull the pin to ground when pressed, so they are active low and need the internal pull-up enabled.
    \item Pins connected to the 8 MHz crystal: PF0 (OSC\_IN, package pin 5) and PF1 (OSC\_OUT, package pin 6), each with a 10 pF stability capacitor to ground.
    \item Default system clock frequency using the internal HSI oscillator: 8 MHz.
    \item Purpose of the loopback header P2: it couples pairs of MCU pins with jumpers (PA8 to PB15, PB14 to PA11, PA7 to PB13) so one peripheral can be wired back into another on the same chip. This allows loopback testing, for example SPI or timer signal verification, without any external equipment.
\end{enumerate}

\section{Task 1: Timer Configuration and Verification}

\subsection{Predictions}

\begin{itemize}
    \item System clock frequency: 8 MHz (HSI, no PLL).
    \item Expected prescaler value: 7999.
    \item Expected ARR value: 999.
    \item Expected interrupt period (s): 1 s.
\end{itemize}

\subsection{Calculations}
The prescaler is chosen to divide the 8 MHz system clock down to a convenient 1 kHz timer clock, so that one timer count corresponds to exactly 1 ms:

\begin{align*}
    \text{Timer Clock} &= \frac{\text{System Clock}}{\text{Prescaler} + 1} = \frac{8\,000\,000}{7999 + 1} = 1000~\text{Hz} \\[2mm]
    \text{ARR} &= \frac{\text{Timer Clock}}{\text{Update Frequency}} - 1 = \frac{1000}{1} - 1 = 999
\end{align*}

Check: period $= \dfrac{(7999+1)(999+1)}{8\,000\,000} = 1$ s.

The ISR toggles PB0 once per interrupt, so the pin produces a square wave with a period of twice the interrupt period. The expected waveform on PB0 is therefore a 0.5 Hz square wave with a 2 s period.

\subsection{Measurement and Verification}
\begin{itemize}
    \item Measured period from the oscilloscope: [MEASURE: toggle period on PB0, expect close to 2 s]
    \item Comparison with your expected period: [MEASURE: state percentage difference; small deviation is expected because the HSI RC oscillator has a factory trimming tolerance of about 1 percent, unlike the crystal]
\end{itemize}

\section{Task 2: LED Control with Timer Interrupt}
Confirm your running light video is in your submission folder, or link it here: [ADD: video link or filename in the GitHub repo]

\section{Task 3: Button Input and Software Debounce}

\subsection{Records}
\begin{itemize}
    \item Expected switch bounce duration: 1 to 20 ms. Small tactile switches typically settle within a few milliseconds, and rarely exceed 10 to 20 ms.
    \item Maximum bounce duration measured on the oscilloscope: [MEASURE: single-shot capture on PA0, worst case over several presses]
    \item Chosen software debounce delay time: 50 ms. This sits well above the measured bounce window, so no bounce edge can retrigger, and well below the fastest intentional repeated press, so no real press is ignored. [ADJUST if your measured bounce exceeds 30 ms]
    \item Formula to calculate the ARR value from the desired period in milliseconds: with the prescaler fixed at 7999 the timer clock is 1 kHz, so one count equals 1 ms and
    \[
        \text{ARR} = \frac{f_{\text{timer}} \cdot T_{\text{ms}}}{1000} - 1 = T_{\text{ms}} - 1.
    \]
    For 1000 ms, ARR $= 999$; for 500 ms, ARR $= 499$.
\end{itemize}

\subsection{Question 4}
The timer waits for the next update event. Our initialisation sets \texttt{AutoReloadPreload = TIM\_AUTORELOAD\_PRELOAD\_ENABLE}, which sets the ARPE bit in TIM16\_CR1. With ARPE set, the ARR register is buffered: a write to TIM16->ARR lands in the preload register, while the counter keeps comparing against the active shadow register. The preload value is only transferred into the shadow register at the next update event, so the cycle in progress completes at the old period and the new frequency takes effect from the following cycle. Resetting CNT to zero restarts the count but does not itself transfer the preload value, since no update event is generated by writing CNT. If ARPE were zero, the write would go straight to the active register and take effect immediately, mid-cycle, which risks a missed compare if the counter is already past the new ARR value.

\section{Task 4: Multi-Mode LED Control}

\subsection{State Machine Explanation}
Mode 3 runs a three-state machine that is polled continuously from the main loop and never blocks, using \texttt{HAL\_GetTick()} deadlines instead of delay loops.

In SPARKLE\_IDLE the code generates a random 8-bit pattern with \texttt{rand() \& 0xFF}, turns on the LEDs whose bits are set, and records their indices in the \texttt{sparkle\_leds\_on} array together with the count. It then computes a deadline \texttt{sparkle\_display\_until} equal to the current tick plus a random hold time of 100 to 1500 ms, and transitions to SPARKLE\_DISPLAY.

In SPARKLE\_DISPLAY the machine simply compares the current tick against the deadline on every pass through the main loop. Nothing else runs in this state, so button handling stays responsive. When the deadline passes, the code schedules the first turn-off time, current tick plus a random 100 to 150 ms, and transitions to SPARKLE\_TURN\_OFF.

In SPARKLE\_TURN\_OFF an index walks through the stored \texttt{sparkle\_leds\_on} array. Each time the current tick reaches the scheduled turn-off time, one LED is switched off and a new random 100 to 150 ms deadline is generated for the next one. When the index reaches the stored count, every LED from the pattern is off and the machine returns to SPARKLE\_IDLE, which immediately starts a new cycle with a fresh pattern.

Because the state machine only ever compares tick deadlines, the PA0 speed toggle and the mode buttons remain debounced and functional at all times. A mode change calls \texttt{set\_mode()}, which clears all LEDs and resets the machine to SPARKLE\_IDLE, so a stale pattern can never survive a mode switch.

\subsection{Measurements}
Every LED turns on at pattern onset and turns off in its slot of the turn-off sequence. The theoretical minimum on-time is the shortest hold plus the shortest first turn-off delay, $100 + 100 = 200$ ms, for the LED switched off first. The theoretical maximum is the longest hold plus eight turn-off slots at 150 ms each, $1500 + 8 \times 150 = 2700$ ms, for the last LED of a full 8-bit pattern.
\begin{itemize}
    \item Measured minimum LED-on time for Mode 3: [MEASURE: scope or slow-motion video over several cycles]
    \item Measured maximum LED-on time for Mode 3: [MEASURE]
\end{itemize}

\section{Task 5: ADC Interrupts and LED Level Meter}

\subsection{Records}
\begin{itemize}
    \item Predicted 12-bit ADC register value, in decimal, for a 2.0 V input:
    \[
        N = \frac{V_{in}}{V_{ref}} \cdot (2^{12} - 1) = \frac{2.0}{3.3} \cdot 4095 = 2481.8 \approx 2482
    \]
    assuming $V_{ref} = V_{DDA} = 3.3$ V.
    \item Measured multimeter voltage at the moment the 5th LED turns on: [MEASURE]
    \item Calculation comparing your measured voltage against the theoretical threshold: our mapping is $k = \left\lfloor \frac{8N + 2048}{4096} \right\rfloor$ (proportional with rounding), so the 5th LED turns on when $8N + 2048 \geq 5 \cdot 4096$, that is $N \geq 2304$. The theoretical threshold voltage is
    \[
        V_{th} = \frac{2304}{4095} \cdot 3.3 = 1.857~\text{V}.
    \]
    [MEASURE: compare your multimeter reading against 1.857 V and give the percentage error]
\end{itemize}

\subsection{Question}
In the STM32F051x datasheet pinout table, the I/O structure for PA5 is listed as TTa: a 3.3 V-tolerant I/O directly connected to the ADC. Unlike FT (five volt tolerant) pins, a TTa pin must not be driven above $V_{DDA} + 0.3$ V, about 3.6 V. Applying 5.0 V forward-biases the internal ESD protection diode from the pin to the supply rail, injecting a large current through the diode into the 3.3 V rail. The absolute maximum injection current (5 mA per pin) is exceeded, which can overheat and permanently destroy the protection diode, damage the analog switch feeding the ADC, corrupt conversions on other channels, and potentially trigger latch-up of the whole device. The damage can be permanent even if the overvoltage is brief.

\section{Task 6: Software PWM via Timer Interrupts}

Output pin: PB5 (LED D5), as configured in the provided Task 6 skeleton.

\subsection{Records}
\begin{itemize}
    \item Calculated Timer Prescaler and ARR values for exactly 10 kHz:
    \[
        f_{int} = \frac{8\,000\,000}{(\text{PSC}+1)(\text{ARR}+1)} = \frac{8\,000\,000}{(7+1)(99+1)} = 10\,000~\text{Hz}
    \]
    so PSC $= 7$ and ARR $= 99$. The ISR counter divides this by 100, giving a 100 Hz PWM period, and drives the pin high for counts 0 to 29, giving 30 of 100 ticks high, a 30 percent duty cycle with an expected pulse width of 3 ms.
    \item Measured total period of the signal: [MEASURE: expect 10 ms]
    \item Measured positive pulse width: [MEASURE: expect 3 ms]
\end{itemize}

\subsection{Question}
A resolution of 1 percent means the period must be divided into 100 timer ticks, so the interrupt must fire 100 times per output period:
\[
    f_{int} = f_{PWM} \cdot \frac{100~\text{percent}}{1~\text{percent}} = 50 \cdot 100 = 5000~\text{Hz} = 5~\text{kHz}.
\]

\end{document}
